% figures09.tex -- Chapter 9 figures, rebuilt 2026-08-16 from the iPad scans.
%
% Every figure below was measured off a scan crop (50/100 px grid overlay,
% tools/ + scratchpad/ch09/grid.py) rather than eyeballed from the photo PDF.
% Constrained points are DERIVED, never placed by hand:
%   parallels  -> equal ratios along the two sides   ($(X)!r!(Y)$)
%   altitudes  -> foot of perpendicular              ($(X)!(P)!(Y)$)
%   crossings  -> (intersection of A--B and C--D)
%   bisectors  -> ratio from the angle-bisector theorem
% so the hypothesis stated in the book holds exactly in the drawing.
% Hypotheses are re-asserted numerically in tools/constraints/ch09.py.
%
% Macro names are chapter-prefixed: \FIGIX<NUMBERWORD>.
%
% Line treatment follows the book's own convention: solid where the book
% prints solid (given), dashed where the book prints dashed (aux).  The
% heavier `key' weight is GONE from this chapter (2026-08-17): Caleb flagged
% Figs 9-23/9-25/9-14/9-27/9-50 for drawing one segment heavier than the rest
% where the book strikes every line at a single weight.  Do not reintroduce it.

\usetikzlibrary{decorations.pathreplacing,arrows.meta}

% The book's fine dashed rule (Fig 9-24's DE, Fig 9-50's compass arcs, the
% dashed base extension of Fig 9-54): thinner and shorter-dashed than `aux'.
\tikzset{
  finedash/.style={fig, line width=0.55pt, dash pattern=on 2.4pt off 1.9pt},
  longdash/.style={fig, dash pattern=on 5.2pt off 3.0pt},
  lead/.style={\clB, line width=0.5pt, -{Latex[length=3.6pt,width=2.6pt]}},
}

% Primed labels.  The math prime sits ~4pt above the letter, which makes the
% collision audit read it as a separate word overlapping its own letter; a
% 1.4pt drop keeps the pair in one word and is visually indistinguishable.
\newcommand{\IXpr}{\raisebox{-1.3pt}{\ensuremath{{}^{\prime}}}}

% Chapter-local label clearance.  brumfiel.sty sets outer sep = 3.4pt, which
% is not quite enough when a labelled point lies ON a line: the line clips the
% corner of the box.  4.8pt clears it without the labels drifting loose.
\tikzset{
  vlab/.append style={outer sep=2.9pt},
  slab/.append style={outer sep=2.9pt},
  klab/.append style={outer sep=2.9pt},
}

% Book-style dimension brace (Figs 9-14, 9-45, 9-46).
\tikzset{
  bkbrace/.style={fig, line width=0.5pt, decorate,
                  decoration={brace, amplitude=3.2pt, raise=1pt}},
  bkbracer/.style={fig, line width=0.5pt, decorate,
                  decoration={brace, mirror, amplitude=3.2pt, raise=1pt}},
}

% ---------------------------------------------------------------- Fig 9-1
% Three parallels l, l', l''; transversal m leans left going down, m' leans
% right; m'' through A and m''' through B are parallel to m' (dashed in the
% book).  Spacing of the parallels is equal because AB = BC on m.
\newcommand{\FIGIXONE}{\begin{bkfigure}[\linewidth]
\begin{tikzpicture}[scale=1.133]
  \foreach \y/\nm in {2.2/$l$, 1.1/$l\IXpr$, 0/$l\IXpr\IXpr$}
    {\draw[given] (0.35,\y) -- (9.05,\y);
     \node[vlab, right] at (9.05,\y) {\nm};}
  % transversal m : dx/dy = -0.527
  \coordinate (A) at (2.90,2.2);
  \coordinate (B) at (2.32,1.1);
  \coordinate (C) at (1.75,0);
  \draw[given] (3.19,2.75) -- (1.57,-0.35);
  % transversal m' : dx/dy = +0.341
  \coordinate (Ap) at (6.90,2.2);
  \coordinate (Bp) at (7.275,1.1);
  \coordinate (Cp) at (7.65,0);
  \draw[given] (6.71,2.75) -- (7.77,-0.35);
  % m'' through A, m''' through B, both parallel to m'
  \coordinate (D) at (3.275,1.1);
  \coordinate (E) at (2.695,0);
  \draw[aux] (A) -- (3.75,-0.30);
  \draw[aux] (B) -- (2.80,-0.30);
  \foreach \p/\n/\d in {A/A/above left, B/B/above left, C/C/above left,
                        D/D/above right, E/E/above right,
                        Ap/A\IXpr/above right, Bp/B\IXpr/above right, Cp/C\IXpr/above right}
    {\node[vlab,\d] at (\p) {$\n$};}
  \node[vlab, below] at (1.57,-0.35) {$m$};
  \node[vlab, below] at (2.80,-0.30) {$m\IXpr\IXpr\IXpr$};
  \node[vlab, below] at (3.75,-0.30) {$m\IXpr\IXpr$};
  \node[vlab, below] at (7.77,-0.35) {$m\IXpr$};
\end{tikzpicture}
\figcap{9--1}
\end{bkfigure}}

% ------------------------------------------------------------ Fig 9-2, 9-3
% Same construction twice: D, E midpoints of AB, BC.  In 9-3 the line m
% through D parallel to AC is carried past E as a dashed ray.
\newcommand{\FIGIXTWOTHREE}{\begin{bkfigure}[\linewidth]
\begin{minipage}{0.36\linewidth}\centering
\begin{tikzpicture}[scale=0.70]
  \coordinate (A) at (0,0); \coordinate (C) at (3.93,0);
  \coordinate (B) at (1.87,1.75);
  \draw[given] (A) -- (C) -- (B) -- cycle;
  \coordinate (D) at ($(A)!0.5!(B)$);
  \coordinate (E) at ($(B)!0.5!(C)$);
  \draw[given] (D) -- (E);
  \foreach \p/\n/\d in {A/A/below left, C/C/below right, B/B/above,
                        D/D/above left, E/E/above right}
    {\node[vlab,\d] at (\p) {$\n$};}
\end{tikzpicture}
\figcap{9--2}
\end{minipage}\hfill
\begin{minipage}{0.47\linewidth}\centering
\begin{tikzpicture}[scale=0.70]
  \coordinate (A) at (0,0); \coordinate (C) at (5.65,0);
  \coordinate (B) at (1.83,1.73);
  \draw[given] (A) -- (C) -- (B) -- cycle;
  \coordinate (D) at ($(A)!0.5!(B)$);
  \coordinate (E) at ($(B)!0.5!(C)$);
  \draw[given] (D) -- (E);
  \draw[aux] (E) -- ($(D)!1.42!(E)$);
  \node[vlab,right] at ($(D)!1.42!(E)$) {$m$};
  \foreach \p/\n/\d in {A/A/below left, C/C/below right, B/B/above,
                        D/D/above left, E/E/above}
    {\node[vlab,\d] at (\p) {$\n$};}
\end{tikzpicture}
\figcap{9--3}
\end{minipage}
\end{bkfigure}}

% ---------------------------------------------------------------- Fig 9-4
% D, E, G are the midpoints of AB, BC, CA.  m' is the line through E parallel
% to AB; it meets AC in G.
\newcommand{\FIGIXFOUR}{\begin{bkfigure}[0.62\linewidth]
\begin{tikzpicture}[scale=1.667]
  \coordinate (A) at (0,0); \coordinate (C) at (2.39,0);
  \coordinate (B) at (0.69,0.98);
  \draw[given] (A) -- (C) -- (B) -- cycle;
  \coordinate (D) at ($(A)!0.5!(B)$);
  \coordinate (E) at ($(B)!0.5!(C)$);
  \coordinate (G) at ($(A)!0.5!(C)$);
  \coordinate (Mp) at ($(G)!1.60!(E)$);
  \draw[given] (D) -- (E);
  % m' is a LINE: the book runs it past G on the far side as well.
  \draw[given] ($(E)!1.34!(G)$) -- (Mp);
  \node[vlab,above right] at (Mp) {$m\IXpr$};
  \foreach \p/\n/\d in {A/A/below left, C/C/right, B/B/above,
                        D/D/above left, G/G/below right}
    {\node[vlab,\d] at (\p) {$\n$};}
  \node[vlab] at ($(E)+(0.41,0.13)$) {$E$};   % right of m', as in the book
\end{tikzpicture}
\figcap{9--4}
\end{bkfigure}}

% ------------------------------------------------------------ Fig 9-5..9-7
% SPLIT 2026-08-17: Ex. 9-2 cites 9-5 in #2, 9-6 in #4/5/6 and 9-7 in #7, so
% each now stands alone at full width and is set inline after its exercise.
\newcommand{\FIGIXFIVE}{\begin{bkfigure}[0.62\linewidth]
\begin{tikzpicture}[scale=1.55]
  % 9-5: triangle EFG, C the midpoint of FG, D fixed by CD || EG and DG || EF,
  % which forces CD = 1/2 EG exactly (the statement to be proved).
  % Vertices measured off the p.156 figure at 400 dpi (E,F,G,C,D all checked:
  % C is the midpoint of FG and D sits at C + (G-E)/2, as the two parallels
  % require).  Do NOT draw to $(C)+(G)-(E)$ -- that is twice too far out and
  % breaks DG || EF.
  \coordinate (E) at (0,0); \coordinate (G) at (3.88,-0.12);
  \coordinate (F) at (0.60,1.338);
  \coordinate (C) at ($(F)!0.5!(G)$);
  \coordinate (D) at ($(C)!0.5!($(C)+(G)-(E)$)$);
  \draw[given] (E) -- (F) -- (G) -- cycle;
  \draw[given] (G) -- (D);
  \draw[given]   (C) -- (D);
  \foreach \p/\n/\d in {E/E/left, F/F/above, G/G/below right}
    {\node[vlab,\d] at (\p) {$\n$};}
  \node[vlab,above] at (C) {$C$};
  \node[vlab,above right] at (D) {$D$};
\end{tikzpicture}
\figcap{9--5}
\end{bkfigure}}

\newcommand{\FIGIXSIX}{\begin{bkfigure}[0.60\linewidth]
\begin{tikzpicture}[scale=1.45]
  % 9-6: three parallels, unequal spacing (the exercises vary AB : BC),
  % cut by transversals m (leaning left) and m' (leaning right).
  % spacing AB : BC = 1.90 : 1, measured off the p.156 figure at 400 dpi.
  \foreach \y in {0, 0.70, 2.03}{\draw[given] (-0.62,\y) -- (3.35,\y);}
  \coordinate (C) at (0,0); \coordinate (B) at (0.0854,0.70);
  \coordinate (A) at (0.2477,2.03);
  \coordinate (Cp) at (2.85,0); \coordinate (Bp) at (2.7716,0.70);
  \coordinate (Ap) at (2.6226,2.03);
  \draw[given] (-0.0915,-0.75) -- (0.3577,2.95);
  \draw[given] (2.934,-0.75)  -- (2.5136,2.95);
  \node[vlab,above] at (0.3577,2.95) {$m$};
  \node[vlab,above] at (2.5136,2.95) {$m\IXpr$};
  \foreach \p/\n/\d in {A/A/above left, B/B/above left, C/C/below left,
                        Ap/A\IXpr/above right, Bp/B\IXpr/above right, Cp/C\IXpr/below right}
    {\node[vlab,\d] at (\p) {$\n$};}
\end{tikzpicture}
\figcap{9--6}
\end{bkfigure}}

\newcommand{\FIGIXSEVEN}{\begin{bkfigure}[0.55\linewidth]
\begin{tikzpicture}[scale=2.30]
  % 9-7: F on EG and P on EQ at the same ratio 2/5, hence FP || GQ.
  \coordinate (G) at (0,0); \coordinate (Q) at (2.35,0);
  \coordinate (E) at (0.80,1.15);
  \draw[given] (G) -- (Q) -- (E) -- cycle;
  \coordinate (F) at ($(E)!0.4!(G)$);
  \coordinate (P) at ($(E)!0.4!(Q)$);
  \draw[given] (F) -- (P);
  \foreach \p/\n/\d in {G/G/below left, Q/Q/below right, E/E/above,
                        F/F/left, P/P/above right}
    {\node[vlab,\d] at (\p) {$\n$};}
\end{tikzpicture}
\figcap{9--7}
\end{bkfigure}}

% ---------------------------------------------------------------- Fig 9-8
% AD = DE = EF = FG on the ray AC; the parallels to BG through D, E, F cut AB
% at its quarter points H, I, J -- so every crossing is derived, not placed.
\newcommand{\FIGIXEIGHT}{\begin{bkfigure}[0.72\linewidth]
\begin{tikzpicture}[scale=1.479]
  \coordinate (A) at (0,0); \coordinate (C) at (5.30,0);
  \coordinate (B) at (3.02,1.50);
  \coordinate (G) at (4.22,0);
  \draw[given] (A) -- (C);
  \draw[given] (A) -- ($(A)!1.09!(B)$);
  \draw[given] (B) -- (G);
  \foreach \i in {1,2,3}{%
     \coordinate (P\i) at ($(A)!\i/4!(G)$);
     \coordinate (H\i) at ($(A)!\i/4!(B)$);
     \draw[given] (H\i) -- (P\i);}
  \node[vlab,left] at (A) {$A$};   \node[vlab,right] at (C) {$C$};
  \node[vlab,above left] at (B) {$B$};
  \node[vlab,below] at (P1) {$D$}; \node[vlab,below] at (P2) {$E$};
  \node[vlab,below] at (P3) {$F$}; \node[vlab,below] at (G) {$G$};
  \node[vlab,above=2.5pt] at (H1) {$H$}; \node[vlab,above=2.5pt] at (H2) {$I$};
  \node[vlab,above=2.5pt] at (H3) {$J$};
\end{tikzpicture}
\figcap{9--8}
\end{bkfigure}}

% ------------------------------------------------------------- Fig 9-9
% SPLIT 2026-08-17 out of \FIGIXNINEFOURTEEN: Ex. 9-3 cites 9-9 in #5/#6,
% 9-10 in #7/#8, 9-11 in #9, 9-12 in #11, 9-13 in #12 and 9-14 in #13, so the
% six now stand alone and are set inline after the exercise that names them.
% 9-9: D on AB, E on BC with BD/BA = BE/BC, hence DE || AC exactly.
\newcommand{\FIGIXNINE}{\begin{bkfigure}[0.70\linewidth]
\begin{tikzpicture}[scale=1.45]
  \coordinate (A) at (0,0); \coordinate (C) at (4.12,0.42);
  \coordinate (B) at (1.32,2.32);
  \draw[given] (A) -- (C) -- (B) -- cycle;
  \coordinate (D) at ($(A)!0.172!(B)$);
  \coordinate (E) at ($(B)!0.828!(C)$);
  \draw[given] (D) -- (E);
  % The book's "Parallel" note is a FAN: single-headed leaders from one apex
  % at the word, one per line named.  (Caleb: our double-headed spans are
  % wrong -- see 9-10, 9-12, 9-14.)
  \coordinate (PA) at (3.44,1.66);
  \node[klab] at (3.30,2.02) {Parallel};
  \draw[lead] (PA) -- ($(D)!0.70!(E)$);
  \draw[lead] (PA) -- ($(A)!0.80!(C)$);
  \foreach \p/\n/\d in {A/A/below left, C/C/right, B/B/above,
                        D/D/left, E/E/above right}
    {\node[vlab,\d] at (\p) {$\n$};}
\end{tikzpicture}
\figcap{9--9}
\end{bkfigure}}

% ------------------------------------------------------------ Fig 9-10
% REBUILT 2026-08-17 against the book photograph (feedback/small/20260817_102341).
% What was wrong: T1 missed A entirely, T1/T2 were drawn near-vertical instead
% of the book's ~31 deg off vertical, the three parallels were dead horizontal
% and cut off just past the outer transversals, and both "Parallel" notes were
% double-headed arrows.  Measured off the photograph: the parallels slope
% 8 deg down to the right, spacing AB : BC = 1.5 : 1, the two slanted
% transversals run 0.586 across per 1 down, and both notes are fans of
% single-headed leaders.
\newcommand{\FIGIXTEN}{\begin{bkfigure}[0.94\linewidth]
\begin{tikzpicture}[scale=1.32]
  % three parallels  y = c - 0.145 x  for c = 2.00, 0.80, 0
  \foreach \c in {2.00,0.80,0}
    {\draw[given] (-0.30,{\c+0.0435}) -- (5.60,{\c-0.812});}
  % transversal 1 through A, nearly vertical, leaning LEFT as it descends
  \coordinate (A) at (0.20,1.971);
  \coordinate (B) at (0.0817,0.7882);
  \coordinate (C) at (0.00286,-0.0004);
  \draw[given] (0.255,2.521) -- (-0.045,-0.479);
  % transversal 2 through the SAME point A (the book crosses them on l1)
  \coordinate (D) at (0.96850,0.65957);
  \coordinate (X) at (1.48083,-0.21472);
  \coordinate (T2a) at (-0.0051,2.321);  \coordinate (T2b) at (1.7352,-0.649);
  \draw[given] (T2a) -- (T2b);
  % transversal 3, parallel to it, through E on l1 and F on l3
  \coordinate (E) at (3.30,1.5215);
  \coordinate (F) at (4.58083,-0.66422);
  \coordinate (T3a) at (3.0949,1.8715); \coordinate (T3b) at (4.8353,-1.0985);
  \draw[given] (T3a) -- (T3b);
  % "Parallel" (top right): three leaders fanning from one apex, one per line
  \coordinate (PA) at (5.35,2.02);
  \node[klab] at (5.02,2.40) {Parallel};
  \draw[lead] (PA) -- (5.15,1.2533);
  \draw[lead] (PA) -- (5.20,0.0460);
  \draw[lead] (PA) -- (5.42,-0.7859);
  % "Parallel" (bottom): one leader each side of the word, to the two slants
  \node[klab] at (3.18,-0.72) {Parallel};
  \draw[lead] (2.52,-0.66) -- ($(T2b)!0.06!(T2a)$);
  \draw[lead] (3.84,-0.78) -- ($(T3b)!0.06!(T3a)$);
  \node[vlab] at ($(A)+(-0.20,-0.22)$) {$A$};
  \node[vlab] at ($(B)+(-0.20,0.22)$) {$B$};
  \node[vlab] at ($(C)+(-0.20,0.22)$) {$C$};
  \node[vlab] at ($(D)+(0.18,0.20)$) {$D$};
  \node[vlab] at ($(E)+(0.18,0.20)$) {$E$};
  \node[vlab] at ($(F)+(0.18,0.20)$) {$F$};
\end{tikzpicture}
\figcap{9--10}
\end{bkfigure}}

% ------------------------------------------------------------ Fig 9-11
% AB : BC : CD = 4 : 3 : 2 down the left side; E and F sit on AG at the
% matching ratios, so BE || CF || DG exactly.  AE is marked 8.
\newcommand{\FIGIXELEVEN}{\begin{bkfigure}[0.68\linewidth]
\begin{tikzpicture}[scale=1.35]
  \coordinate (A) at (0,3.15); \coordinate (D) at (0,0);
  \coordinate (G) at (4.55,0);
  \coordinate (B) at (0,1.75); \coordinate (C) at (0,0.70);
  \coordinate (E) at ($(A)!4/9!(G)$);
  \coordinate (F) at ($(A)!7/9!(G)$);
  \draw[given] (A) -- (D) -- (G) -- cycle;
  \draw[given] (B) -- (E); \draw[given] (C) -- (F);
  \node[slab,left=0pt]  at (0,2.45) {4};
  \node[slab,left=0pt]  at (0,1.225) {3};
  \node[slab,left=0pt]  at (0,0.35) {2};
  \node[slab,above] at (0.95,2.60) {8};
  \foreach \p/\n/\d in {A/A/above, D/D/below left, G/G/below right}
    {\node[vlab,\d] at (\p) {$\n$};}
  \node[vlab,left=7pt] at (B) {$B$};
  \node[vlab,left=7pt] at (C) {$C$};
  \node[vlab,above right] at (E) {$E$};
  \node[vlab,above right] at (F) {$F$};
\end{tikzpicture}
\figcap{9--11}
\end{bkfigure}}

% ------------------------------------------------------------ Fig 9-12
% parallels l1,l2,l3 with gaps 2 : 3; a zigzag of three transversals.
% 2 and 3 are the pieces of the first, 4 the lower piece of the second,
% 6 the lower piece of the third; AB is the piece asked for.
% The numerals here are SEGMENT LENGTHS, not angle names -- the book strikes
% no arcs at them and neither do we (checked on the p.158 plate at 500 dpi).
% The "Parallel" note is the book's three-leader fan, not a double-headed span.
\newcommand{\FIGIXTWELVE}{\begin{bkfigure}[0.90\linewidth]
\begin{tikzpicture}[scale=1.15]
  \coordinate (T)  at (0,0);           % foot of transversal 1 on l3
  \coordinate (V1) at (1.189,1.886);   % apex, on l1
  \coordinate (V2) at (3.156,-0.303);  % on l3
  \coordinate (B)  at (6.411,1.3845);  % on l1
  \coordinate (P)  at ($(T)!0.6!(V1)$);
  \coordinate (Q)  at ($(V1)!0.4!(V2)$);
  \coordinate (A)  at ($(V2)!0.6!(B)$);
  \draw[given] (V1) -- (6.62,1.365);   % l1
  \draw[given] (P)  -- (6.60,0.566);   % l2
  \draw[given] (T)  -- (6.62,-0.635);  % l3
  \draw[given] (T) -- (V1) -- (V2) -- (B);
  \coordinate (PA) at (6.85,2.05);
  \node[klab] at (6.58,2.42) {Parallel};
  \draw[lead] (PA) -- (6.45,1.3813);
  \draw[lead] (PA) -- (6.45,0.5805);
  \draw[lead] (PA) -- (6.50,-0.6235);
  \node[slab,left=1pt]  at ($(P)!0.5!(V1)$) {2};
  \node[slab,left]  at ($(T)!0.5!(P)$)  {3};
  \node[slab] at ($($(Q)!0.5!(V2)$)+(0.17,0.15)$) {4};
  \node[slab] at ($(V2)!0.5!(A)+(-0.18,0.30)$) {6};
  \node[vlab] at ($(A)+(-0.22,0.28)$) {$A$};
  \node[vlab,above left] at (B) {$B$};
\end{tikzpicture}
\figcap{9--12}
\end{bkfigure}}

% ------------------------------------------------------------ Fig 9-13
% parallelogram ABCD, E and F the midpoints of AB and CD; the only diagonal
% drawn is AC, cut by ED and BF at its trisection points G and H.
% vertices measured off the p.158 figure at 400 dpi: AD runs 1.5 deg BELOW
% the horizontal and AB stands at 77 deg, giving the book's upright shape.
\newcommand{\FIGIXTHIRTEEN}{\begin{bkfigure}[0.66\linewidth]
\begin{tikzpicture}[scale=1.40]
  \coordinate (A) at (0,0); \coordinate (D) at (3.70,-0.10);
  \coordinate (B) at (0.60,2.55); \coordinate (C) at (4.30,2.45);
  \draw[given] (A) -- (D) -- (C) -- (B) -- cycle;
  \coordinate (E) at ($(A)!0.5!(B)$);
  \coordinate (F) at ($(C)!0.5!(D)$);
  \draw[given] (A) -- (C);
  \draw[given] (E) -- (D); \draw[given] (B) -- (F);
  \coordinate (G) at (intersection of A--C and E--D);
  \coordinate (H) at (intersection of A--C and B--F);
  \foreach \p/\n/\d in {A/A/below left, D/D/below right, C/C/above right,
                        B/B/above, E/E/left, F/F/right}
    {\node[vlab,\d] at (\p) {$\n$};}
  \node[vlab] at ($(G)+(0,0.42)$) {$G$};
  \node[vlab] at ($(H)+(0,0.42)$) {$H$};
\end{tikzpicture}
\figcap{9--13}
\end{bkfigure}}

% ------------------------------------------------------------ Fig 9-14
% D on LT and E on LR at the same ratio, so DE || TR; the braces mark
% x = LT, 1 = LD, y = LE, z = LR.  x : 1 = z : y is the point of the problem.
% REBUILT 2026-08-17 from the book photograph (feedback/small/20260817_102437):
% the two brace rows on each side were nearly on top of the side they measure,
% so the nesting could not be read, and the x brace was drawn INSIDE the 1
% brace (the book has the long brace outermost on both sides).  Offsets are
% now measured off the photograph: 0.07 and 0.24 of the side's own length.
% The "Parallel" leaders now fan from ONE apex at the word, as the book draws
% them, and the upper leader crosses the right side on its way in (so does
% the book's).
\newcommand{\FIGIXFOURTEEN}{\begin{bkfigure}[0.88\linewidth]
\begin{tikzpicture}[scale=1.30]
  \coordinate (L) at (0,0); \coordinate (R) at (4.66,0);
  \coordinate (T) at (2.16,2.12);
  \coordinate (D) at ($(L)!0.5!(T)$);
  \coordinate (E) at ($(L)!0.5!(R)$);
  \draw[given] (L) -- (R) -- (T) -- cycle;
  \draw[given] (D) -- (E);
  % outward unit normal of LT is (-0.7003, 0.7136); of LR it is (0,-1)
  % Drawn L -> T (not T -> L) so the brace teeth point AWAY from the side, out
  % towards their labels -- which is the way the base braces below already run,
  % and the way the book strikes all four.
  \draw[bkbrace] ($(L)+(-0.5042,0.5138)$) -- ($(T)+(-0.5042,0.5138)$);
  \draw[bkbrace] ($(L)+(-0.1541,0.1570)$) -- ($(D)+(-0.1541,0.1570)$);
  \draw[bkbracer] ($(L)+(0,-0.24)$) -- ($(E)+(0,-0.24)$);
  \draw[bkbracer] ($(L)+(0,-0.76)$) -- ($(R)+(0,-0.76)$);
  \node[slab] at (0.422,1.731) {$x$};
  \node[slab] at (0.232,0.844) {1};
  \node[slab] at (1.165,-0.52) {$y$};
  \node[slab] at (2.330,-1.04) {$z$};
  \coordinate (PA) at (3.62,1.28);
  \node[klab,right] at (3.74,1.30) {Parallel};
  \draw[lead] (PA) -- ($(D)!0.5!(E)$);
  \draw[lead] (PA) -- ($(T)!0.55!(R)$);
\end{tikzpicture}
\figcap{9--14}
\end{bkfigure}}

% --------------------------------------------------------------- Fig 9-15
% Sides 5, 6, 7 built exactly; the primed triangle is the same shape at 2/5.
\newcommand{\FIGIXFIFTEEN}{\begin{bkfigure}[0.86\linewidth]
\begin{tikzpicture}[scale=1.160]
  \coordinate (E) at (0,0); \coordinate (G) at (3.5,0);
  \coordinate (F) at (1.357,2.099);
  \draw[given] (E) -- (G) -- (F) -- cycle;
  \node[slab,left]  at ($(E)!0.5!(F)$) {5};
  \node[slab] at ($($(F)!0.5!(G)$)+(0.154,0.157)$) {6};
  \node[slab,below] at ($(E)!0.5!(G)$) {7};
  \node[vlab,below left] at (E) {$E$}; \node[vlab,below right] at (G) {$G$};
  \node[vlab,above] at (F) {$F$};
  \begin{scope}[xshift=5.1cm, yshift=0.02cm]
    \coordinate (Ep) at (0,0); \coordinate (Gp) at (1.4,0);
    \coordinate (Fp) at (0.543,0.840);
    \draw[given] (Ep) -- (Gp) -- (Fp) -- cycle;
    \node[slab,left] at ($(Ep)!0.5!(Fp)$) {2};
    \node[vlab,below left] at (Ep) {$E\IXpr$}; \node[vlab,below right] at (Gp) {$G\IXpr$};
    \node[vlab,above] at (Fp) {$F\IXpr$};
  \end{scope}
\end{tikzpicture}
\figcap{9--15}
\end{bkfigure}}

% --------------------------------------------------------------- Fig 9-16
\newcommand{\FIGIXSIXTEEN}{\begin{bkfigure}[0.9\linewidth]
\begin{tikzpicture}[scale=1.377]
  \draw[given] (0,0.62)   -- (2.75,0.62);
  \draw[given] (0.10,0)   -- (2.40,0);
  \draw[given] (4.35,0.62) -- (5.50,0.62);
  \draw[given] (4.65,0)   -- (5.25,0);
  \foreach \x/\y in {0/0.62, 2.75/0.62, 0.10/0, 2.40/0,
                     4.35/0.62, 5.50/0.62, 4.65/0, 5.25/0}
    {\dt{\x,\y}}
  \node[vlab,left]  at (0,0.62)    {$A$};
  \node[vlab,right] at (2.75,0.62) {$B$};
  \node[vlab,left]  at (0.10,0)    {$A\IXpr$};
  \node[vlab,right] at (2.40,0)    {$B\IXpr$};
  \node[vlab,left]  at (4.35,0.62) {$C$};
  \node[vlab,right] at (5.50,0.62) {$D$};
  \node[vlab,left]  at (4.65,0)    {$E$};
  \node[vlab,right] at (5.25,0)    {$F$};
\end{tikzpicture}
\figcap{9--16}
\end{bkfigure}}

% --------------------------------------------------------------- Fig 9-17
% CD cut into four congruent subintervals; EF is one of them.
\newcommand{\FIGIXSEVENTEEN}{\begin{bkfigure}[0.6\linewidth]
\begin{tikzpicture}[scale=1.958]
  \draw[given] (0,0) -- (2.28,0);
  \foreach \x in {0,0.57,1.14,1.71,2.28}{\dt{\x,0}}
  \node[vlab,left]  at (0,0)    {$C$};
  \node[vlab,right] at (2.28,0) {$D$};
  \node[vlab,above=2.6pt] at (0.57,0) {$E$};
  \node[vlab,above=2.6pt] at (1.14,0) {$F$};
\end{tikzpicture}
\figcap{9--17}
\end{bkfigure}}

% --------------------------------------------------------- Fig 9-18, 9-19
\newcommand{\FIGIXEIGHTEENNINETEEN}{\begin{bkfigure}[\linewidth]
\begin{minipage}{0.52\linewidth}\centering
\begin{tikzpicture}[scale=0.55]
  \draw[given] (0,0) -- (3.6,0) -- (1.6,2.2) -- cycle;
  \node[vlab,below left] at (0,0) {$A$}; \node[vlab,below right] at (3.6,0) {$C$};
  \node[vlab,above] at (1.6,2.2) {$B$};
  \begin{scope}[xshift=6.2cm, scale=0.55]
    \draw[given] (0,0) -- (3.6,0) -- (1.0,2.28) -- cycle;
    \node[vlab,below left] at (0,0) {$A\IXpr$}; \node[vlab,below right] at (3.6,0) {$C\IXpr$};
    \node[vlab,above] at (1.6,2.2) {$B\IXpr$};
  \end{scope}
\end{tikzpicture}
\figcap{9--18}
\end{minipage}\hfill
\begin{minipage}{0.42\linewidth}\centering
\begin{tikzpicture}[scale=0.57]
  % 9-19: convex pentagon ABCDE.  Vertices measured off the p.163 figure at
  % 400 dpi.  The cycle is A-B-C-D-E (an earlier draft joined E to B and D to
  % A, i.e. it swapped A and E in the path, which put A almost on top of E).
  \def\pentA{(0,1.38) (0.66,2.47) (2.25,1.67) (1.87,0.31) (0.05,0)}
  \draw[given] plot coordinates {\pentA} -- cycle;
  \node[vlab,left]  at (0,1.38)    {$A$};
  \node[vlab,above] at (0.66,2.47) {$B$};
  \node[vlab,right] at (2.25,1.67) {$C$};
  \node[vlab,below right] at (1.87,0.31) {$D$};
  \node[vlab,below left]  at (0.05,0)    {$E$};
  \begin{scope}[xshift=4.15cm, yshift=0.30cm, scale=0.62]
    \draw[given] plot coordinates {\pentA} -- cycle;
    \node[vlab,left]  at (0,1.38)    {$A\IXpr$};
    \node[vlab,above] at (0.66,2.47) {$B\IXpr$};
    \node[vlab,right] at (2.25,1.67) {$C\IXpr$};
    \node[vlab,below right] at (1.87,0.31) {$D\IXpr$};
    \node[vlab,below left]  at (0.05,0)    {$E\IXpr$};
  \end{scope}
\end{tikzpicture}
\figcap{9--19}
\end{minipage}
\end{bkfigure}}

% --------------------------------------------------------- Fig 9-20, 9-21
% SPLIT 2026-08-17: Ex. 9-8 cites 9-20 in #15 and 9-21 in #16.
\newcommand{\FIGIXTWENTY}{\begin{bkfigure}[0.54\linewidth]
\begin{tikzpicture}[scale=2.00]
  % 9-20: E, F midpoints of AB and BC.
  \coordinate (A) at (0,0); \coordinate (C) at (2.55,0);
  \coordinate (B) at (1.68,1.10);
  \draw[given] (A) -- (C) -- (B) -- cycle;
  \coordinate (E) at ($(A)!0.5!(B)$);
  \coordinate (F) at ($(B)!0.5!(C)$);
  \draw[given] (E) -- (F);
  \foreach \p/\n/\d in {A/A/below left, C/C/below right, B/B/above,
                        E/E/above left, F/F/above right}
    {\node[vlab,\d] at (\p) {$\n$};}
\end{tikzpicture}
\figcap{9--20}
\end{bkfigure}}

\newcommand{\FIGIXTWENTYONE}{\begin{bkfigure}[0.62\linewidth]
\begin{tikzpicture}[scale=1.95]
  % 9-21: A', B', C', D' are the midpoints of OA, OB, OC, OD.
  % Spoke lengths and directions measured off the p.164 figure at 400 dpi:
  % OA : OB : OC : OD = 428 : 190 : 275 : 200 px.  (An earlier draft had OD
  % 30% longer than OB; the book draws them nearly equal.)
  \coordinate (O) at (0,0);
  \coordinate (A) at (-1.72,0.01);  \coordinate (B) at (-0.05,0.76);
  \coordinate (C) at (1.105,0.07);  \coordinate (D) at (0.30,-0.75);
  \draw[given] (A) -- (B) -- (C) -- (D) -- cycle;
  \foreach \p in {A,B,C,D}{\draw[given] (O) -- (\p);
                           \coordinate (\p p) at ($(O)!0.5!(\p)$);}
  \draw[given] (Ap) -- (Bp) -- (Cp) -- (Dp) -- cycle;
  \node[vlab,left]  at (A) {$A$};   \node[vlab,above] at (B) {$B$};
  \node[vlab,right] at (C) {$C$};   \node[vlab,below] at (D) {$D$};
  \node[vlab] at ($(Ap)+(-0.08,0.19)$)  {$A\IXpr$};
  \node[vlab] at ($(Bp)+(-0.24,0.08)$)  {$B\IXpr$};
  \node[vlab] at ($(Cp)+(0.11,0.115)$)  {$C\IXpr$};
  \node[vlab] at ($(Dp)+(-0.22,-0.10)$) {$D\IXpr$};
  \node[vlab] at ($(O)+(0.13,0.14)$)    {$O$};
\end{tikzpicture}
\figcap{9--21}
\end{bkfigure}}

% --------------------------------------------------------------- Fig 9-22
% AB is cut into five unit pieces by F, A'', E, D; the parallels to AC through
% them meet BC in G, C'', H, I.  The primed triangle is the same shape at 3/5.
\newcommand{\FIGIXTWENTYTWO}{\begin{bkfigure}[0.9\linewidth]
\begin{tikzpicture}[scale=1.885]
  \coordinate (A) at (0,0); \coordinate (C) at (1.75,0);
  \coordinate (B) at (1.02,1.87);
  \draw[given] (A) -- (C) -- (B) -- cycle;
  \coordinate (F)   at ($(A)!1/5!(B)$);   \coordinate (G)   at ($(B)!4/5!(C)$);
  \coordinate (App) at ($(A)!2/5!(B)$);   \coordinate (Cpp) at ($(B)!3/5!(C)$);
  \coordinate (E)   at ($(A)!3/5!(B)$);   \coordinate (H)   at ($(B)!2/5!(C)$);
  \coordinate (D)   at ($(A)!4/5!(B)$);   \coordinate (I)   at ($(B)!1/5!(C)$);
  \draw[aux] (F) -- (G);
  \draw[given] (App) -- (Cpp);
  \draw[aux] (E) -- (H);
  \draw[aux] (D) -- (I);
  \node[vlab,left] at (F) {$F$};   \node[vlab,left] at (App) {$A\IXpr\IXpr$};
  \node[vlab,left] at (E) {$E$};   \node[vlab,left] at (D) {$D$};
  \node[vlab,right] at (G) {$G$};  \node[vlab,right] at (Cpp) {$C\IXpr\IXpr$};
  \node[vlab,right] at (H) {$H$};  \node[vlab,right] at (I) {$I$};
  \node[vlab,below left] at (A) {$A$}; \node[vlab,below right] at (C) {$C$};
  \node[vlab,above] at (B) {$B$};
  \begin{scope}[xshift=2.70cm, scale=0.6]
    \draw[given] (0,0) -- (1.75,0) -- (1.02,1.87) -- cycle;
    \node[vlab,below left] at (0,0) {$A\IXpr$};
    \node[vlab,below right] at (1.75,0) {$C\IXpr$};
    \node[vlab,above] at (1.02,1.87) {$B\IXpr$};
  \end{scope}
\end{tikzpicture}
\figcap{9--22}
\end{bkfigure}}

% ---------------------------------------------------------- Fig 9-23..9-29
% SPLIT 2026-08-17: Ex. 9-9 cites 9-23 in #1, 9-24 in #4, 9-25 in #5, 9-26 in
% #7, 9-27 in #10, 9-28 in #11 and 9-29 in #12; the book sets each beside the
% exercise that names it, so each now stands alone and is set inline.
%
% 9-23: AE = 4, ED = 6, so B and E sit at 0.4 along AC and AD; BE || CD.
% The book runs CD almost horizontally right and SLIGHTLY DOWN to D (the
% render had it dead level), puts E level with its own point rather than down
% by D, and strikes BE at the same weight as the sides.
\newcommand{\FIGIXTWENTYTHREE}{\begin{bkfigure}[0.56\linewidth]
\begin{tikzpicture}[scale=2.00]
  \coordinate (C) at (0,0); \coordinate (D) at (2.10,-0.10);
  \coordinate (A) at (2.58,1.32);
  \draw[given] (C) -- (D) -- (A) -- cycle;
  \coordinate (B) at ($(A)!0.4!(C)$);
  \coordinate (E) at ($(A)!0.4!(D)$);
  \draw[given] (B) -- (E);
  \foreach \p/\n/\d in {C/C/left, D/D/below, A/A/above,
                        B/B/above left, E/E/right}
    {\node[vlab,\d] at (\p) {$\n$};}
\end{tikzpicture}
\figcap{9--23}
\end{bkfigure}}

% 9-24: AB and DE both run north-south; D at 0.85 along AC, E directly below.
% The lake is traced from the iPad scan (scan12 p8) and placed against the
% triangle by matching its bounding box to the printed figure.
% 2026-08-17, against the p.166 plate at 500 dpi:
%   * the lake carries THREE contours, not four -- the innermost oval went;
%   * DE is a FINE DASHED segment in the book, not a heavy solid one;
%   * C was not actually ON the base line (the base ran from B to (3.265,0.075)
%     while C sat at y = 0.025), which is the "kink near B" Caleb saw.  The
%     base is now dead horizontal and AB dead vertical, as "run north-south"
%     requires, and C lies on it exactly.
\newcommand{\FIGIXTWENTYFOUR}{\begin{bkfigure}[0.62\linewidth]
\begin{tikzpicture}[scale=1.75]
  \coordinate (B) at (0,0); \coordinate (A) at (0,2.139);
  \coordinate (C) at (2.377,0);
  \coordinate (D) at ($(A)!0.85!(C)$);
  \coordinate (E) at (D |- B);
  \draw[given, line width=0.5pt]
    plot[smooth cycle, tension=0.7] coordinates {
      (1.636,2.336) (2.026,2.242) (2.307,2.055) (2.463,1.868) (2.588,1.727)
      (2.666,1.540) (2.791,1.384) (2.806,1.197) (2.775,0.994) (2.650,0.791)
      (2.494,0.635) (2.260,0.448) (2.026,0.339) (1.714,0.245) (1.480,0.191)
      (1.246,0.229) (1.012,0.339) (0.841,0.526) (0.669,0.822) (0.466,0.947)
      (0.334,1.150) (0.303,1.353) (0.334,1.509) (0.388,1.727) (0.513,1.914)
      (0.700,2.070) (0.934,2.195) (1.246,2.289)};
  \draw[given, line width=0.5pt]
    plot[smooth cycle, tension=0.7] coordinates {
      (1.558,2.070) (1.839,2.016) (2.058,1.922) (2.214,1.789) (2.292,1.634)
      (2.284,1.478) (2.370,1.337) (2.370,1.150) (2.292,0.978) (2.167,0.838)
      (1.995,0.729) (1.792,0.651) (1.574,0.604) (1.356,0.588) (1.153,0.620)
      (0.966,0.697) (0.810,0.822) (0.700,0.978) (0.622,1.150) (0.576,1.321)
      (0.544,1.493) (0.576,1.665) (0.669,1.821) (0.825,1.946) (1.028,2.024)
      (1.293,2.062)};
  \draw[given, line width=0.5pt]
    plot[smooth cycle, tension=0.7] coordinates {
      (1.558,1.556) (1.761,1.509) (1.855,1.384) (1.855,1.228) (1.792,1.087)
      (1.636,0.978) (1.480,0.900) (1.309,0.838) (1.137,0.830) (1.012,0.900)
      (0.934,1.010) (0.888,1.150) (0.864,1.290) (0.903,1.399) (1.012,1.478)
      (1.184,1.524) (1.371,1.553)};
  \draw[given] (A) -- (B) -- (3.30,0);
  \draw[given] (A) -- (C);
  \draw[finedash] (D) -- (E);
  \node[vlab,above left]  at (A) {$A$};
  \node[vlab,below left]  at (B) {$B$};
  \node[vlab,below] at ($(C)+(0.10,0)$) {$C$};
  % D's label sits in the clear band between contour 1 and contour 2, on the
  % same side as the book's.
  \node[vlab] at ($(D)+(0.07,0.27)$) {$D$};
  \node[vlab,below] at (E) {$E$};
\end{tikzpicture}
\figcap{9--24}
\end{bkfigure}}

% 9-25: D on BA and E on BC at the same ratio, so DE || AC.  Uniform weight:
% the book does NOT strike DE heavier than the sides (Caleb, 2026-08-17).
\newcommand{\FIGIXTWENTYFIVE}{\begin{bkfigure}[0.55\linewidth]
\begin{tikzpicture}[scale=2.10]
  \coordinate (A) at (0,0); \coordinate (C) at (2.40,0);
  \coordinate (B) at (1.30,1.42);
  \draw[given] (A) -- (C) -- (B) -- cycle;
  \coordinate (D) at ($(B)!0.53!(A)$);
  \coordinate (E) at ($(B)!0.53!(C)$);
  \draw[given] (D) -- (E);
  \foreach \p/\n/\d in {A/A/below left, C/C/right, B/B/above}
    {\node[vlab,\d] at (\p) {$\n$};}
  \node[vlab] at ($(D)+(-0.147,0.135)$) {$D$};
  \node[vlab] at ($(E)+(0.158,0.122)$)  {$E$};
\end{tikzpicture}
\figcap{9--25}
\end{bkfigure}}

% 9-26: AB _|_ AE at A and DE _|_ AE at E, on opposite sides, so BD crosses AE
% at C with AC/CE = AB/DE.  2026-08-17: AE now LEANS as the book draws it (the
% render had it dead vertical), DE is the short side and AB the long one in the
% book's 2.5 : 1 proportion, and C's label is lifted clear of the crossing --
% it was sitting on the lines.
\newcommand{\FIGIXTWENTYSIX}{\begin{bkfigure}[0.44\linewidth]
\begin{tikzpicture}[scale=2.80]
  \coordinate (A) at (0.18,1.60); \coordinate (E) at (0,0);
  \coordinate (B) at (1.0346,1.5038); \coordinate (D) at (-0.3379,0.0381);
  \draw[given] (A) -- (E); \draw[given] (A) -- (B); \draw[given] (E) -- (D);
  \draw[given] (B) -- (D);
  \coordinate (C) at (intersection of A--E and B--D);
  \foreach \p/\n/\d in {A/A/above left, B/B/above right, E/E/below right,
                        D/D/below left}
    {\node[vlab,\d] at (\p) {$\n$};}
  \node[vlab] at ($(C)+(-0.20,0.02)$) {$C$};
\end{tikzpicture}
\figcap{9--26}
\end{bkfigure}}

% 9-27: l and m are parallel (both of slope 2); EI and FH cross at G, so
% triangle EFG ~ triangle IHG.  2026-08-17: both parallels are drawn
% appreciably longer (they were stubs, so l || m could not be read), and the
% l/m labels now sit AT the ends of their own lines instead of floating past
% them.  The transversals are no longer heavier than the parallels.
\newcommand{\FIGIXTWENTYSEVEN}{\begin{bkfigure}[0.60\linewidth]
\begin{tikzpicture}[scale=1.70]
  \coordinate (F) at (0,0);    \coordinate (E) at (0.40,0.80);
  \coordinate (I) at (1.98,0); \coordinate (H) at (2.454,0.948);
  \draw[given] ($(F)!-0.70!(E)$) -- ($(F)!1.85!(E)$);
  \draw[given] ($(I)!-0.70!(H)$) -- ($(I)!1.80!(H)$);
  \draw[given] (E) -- (I); \draw[given] (F) -- (H);
  \coordinate (G) at (intersection of E--I and F--H);
  \node[vlab,below left]  at ($(F)!-0.70!(E)$) {$l$};
  \node[vlab,above right] at ($(I)!1.80!(H)$)  {$m$};
  \foreach \p/\n/\d in {I/I/below right, H/H/right, G/G/above}
    {\node[vlab,\d] at (\p) {$\n$};}
  \node[vlab] at ($(E)+(-0.42,0.12)$) {$E$};
  \node[vlab,left]       at (F) {$F$};
\end{tikzpicture}
\figcap{9--27}
\end{bkfigure}}

% 9-28: PS bisects angle P, so QS/SR = PQ/PR (ratio 0.468 for these sides).
% 2026-08-17 -- the figure Caleb called the worst offender.  The R label of the
% left triangle and the P' label of the right one printed jammed together; the
% book leaves a clear horizontal gap between the two triangles.  The gap is now
% 0.55 of a left-triangle unit (the book's is 0.17 of its WIDTH; ours needs
% more because our lettering is proportionally larger), the primed triangle is
% at the book's 0.62 rather than 0.583, and the whole block is drawn at nearly
% five times the old area.
\newcommand{\FIGIXTWENTYEIGHT}{\begin{bkfigure}[0.76\linewidth]
\begin{tikzpicture}[scale=1.90]
  \coordinate (P) at (0,0); \coordinate (R) at (1.80,0);
  \coordinate (Q) at (0.96,1.26);
  \draw[given] (P) -- (R) -- (Q) -- cycle;
  \coordinate (S) at ($(Q)!0.468!(R)$);
  \draw[given] (P) -- (S);
  \foreach \p/\n/\d in {P/P/below left, R/R/below right, Q/Q/above, S/S/right}
    {\node[vlab,\d] at (\p) {$\n$};}
  % primed triangle: same shape at 0.62, set 0.55 units clear of R
  \coordinate (Pp) at (2.35,0); \coordinate (Rp) at (3.466,0);
  \coordinate (Qp) at (2.9452,0.7812);
  \draw[given] (Pp) -- (Rp) -- (Qp) -- cycle;
  \coordinate (Sp) at ($(Qp)!0.468!(Rp)$);
  \draw[given] (Pp) -- (Sp);
  \foreach \p/\n/\d in {Pp/P\IXpr/below left, Rp/R\IXpr/below right,
                        Qp/Q\IXpr/above, Sp/S\IXpr/right}
    {\node[vlab,\d] at (\p) {$\n$};}
\end{tikzpicture}
\figcap{9--28}
\end{bkfigure}}

% 9-29: CD and AE are altitudes -- D and E are true perpendicular feet.
% The book strikes no right-angle squares at D or E; none are drawn.
\newcommand{\FIGIXTWENTYNINE}{\begin{bkfigure}[0.54\linewidth]
\begin{tikzpicture}[scale=2.60]
  \coordinate (A) at (0,0); \coordinate (C) at (1.88,0);
  \coordinate (B) at (1.15,1.45);
  \draw[given] (A) -- (C) -- (B) -- cycle;
  \coordinate (D) at ($(A)!(C)!(B)$);
  \coordinate (E) at ($(B)!(A)!(C)$);
  \draw[given] (C) -- (D); \draw[given] (A) -- (E);
  \coordinate (F) at (intersection of C--D and A--E);
  \foreach \p/\n/\d in {A/A/below left, C/C/below right, B/B/above,
                        D/D/above left, E/E/above right}
    {\node[vlab,\d] at (\p) {$\n$};}
  \node[vlab] at ($(F)+(0.00,0.36)$) {$F$};
\end{tikzpicture}
\figcap{9--29}
\end{bkfigure}}

% --------------------------------------------------------------- Fig 9-30
% AD/AB = AE/AC, so DE || BC.  DF is the auxiliary line of the proof, drawn
% to a point F beyond E on AC -- the proof's job is to show F = E.
\newcommand{\FIGIXTHIRTY}{\begin{bkfigure}[0.52\linewidth]
\begin{tikzpicture}[scale=1.522]
  \coordinate (B) at (0,0); \coordinate (C) at (2.30,0.20);
  \coordinate (A) at (1.20,1.30);
  \draw[given] (B) -- (C) -- (A) -- cycle;
  \coordinate (D) at ($(A)!0.50!(B)$);
  \coordinate (E) at ($(A)!0.50!(C)$);
  \coordinate (F) at ($(A)!0.68!(C)$);
  \draw[given] (D) -- (E);
  \draw[aux] (D) -- (F);
  \node[vlab] at ($(D)+(-0.162,0.149)$) {$D$};
  \foreach \p/\n/\d in {B/B/below left, C/C/right, A/A/above,
                        E/E/above right}
    {\node[vlab,\d] at (\p) {$\n$};}
  \node[vlab] at ($(F)+(0.36,0.03)$) {$F$};
\end{tikzpicture}
\figcap{9--30}
\end{bkfigure}}

% --------------------------------------------------------------- Fig 9-31
% B'' on AB and C'' on AC at the same ratio, so B''C'' || BC (Statement 4).
\newcommand{\FIGIXTHIRTYONE}{\begin{bkfigure}[0.84\linewidth]
\begin{tikzpicture}[scale=1.667]
  \coordinate (A) at (0,0); \coordinate (C) at (1.63,0);
  \coordinate (B) at (1.03,1.13);
  \draw[given] (A) -- (C) -- (B) -- cycle;
  \coordinate (Bpp) at ($(A)!0.46!(B)$);
  \coordinate (Cpp) at ($(A)!0.46!(C)$);
  \draw[given] (Bpp) -- (Cpp);
  \foreach \p/\n/\d in {A/A/left, C/C/right, B/B/above,
                        Bpp/B\IXpr\IXpr/above left, Cpp/C\IXpr\IXpr/below}
    {\node[vlab,\d] at (\p) {$\n$};}
  \begin{scope}[xshift=2.85cm, scale=0.54]
    \draw[given] (0,0) -- (1.63,0) -- (1.03,1.13) -- cycle;
    \node[vlab,below left] at (0,0) {$A\IXpr$};
    \node[vlab,below right] at (1.63,0) {$C\IXpr$};
    \node[vlab,above] at (1.03,1.13) {$B\IXpr$};
  \end{scope}
\end{tikzpicture}
\figcap{9--31}
\end{bkfigure}}

% ---------------------------------------------------------- Fig 9-32..9-35
% SPLIT 2026-08-17: Ex. 9-10 cites 9-32 in #2, 9-33 in #4, 9-34 in #5 and
% 9-35 in *#6, so each stands alone and is set inline after its exercise.
\newcommand{\FIGIXTHIRTYTWO}{\begin{bkfigure}[0.58\linewidth]
\begin{tikzpicture}[scale=1.55]
  % 9-32: O divides both AD and BC in the ratio 2 : 1, matching AO = 10,
  % OD = 5, BO = 12, OC = 6.
  % AB : AD = 0.27 in the book (measured p.169 at 400 dpi); an earlier draft
  % drew it at 0.44, which made the figure far too tall.
  \coordinate (O) at (1.35,0);
  \coordinate (A) at (0,-0.06);  \coordinate (B) at (-0.03,-0.60);
  \coordinate (D) at ($(A)!1.5!(O)$);
  \coordinate (C) at ($(B)!1.5!(O)$);
  \draw[given] (A) -- (B); \draw[given] (C) -- (D);
  \draw[given] (A) -- (D); \draw[given] (B) -- (C);
  \foreach \p/\n/\d in {A/A/above left, B/B/below left, C/C/above right,
                        D/D/right, O/O/above}
    {\node[vlab,\d] at (\p) {$\n$};}
\end{tikzpicture}
\figcap{9--32}
\end{bkfigure}}

\newcommand{\FIGIXTHIRTYTHREE}{\begin{bkfigure}[0.55\linewidth]
\begin{tikzpicture}[scale=2.00]
  % 9-33: AE = x.AC and AD = x.AB with the same x, so ED || BC.
  \coordinate (B) at (0,0); \coordinate (C) at (2.50,0);
  \coordinate (A) at (1.03,1.44);
  \draw[given] (B) -- (C) -- (A) -- cycle;
  \coordinate (E) at ($(A)!0.44!(B)$);
  \coordinate (D) at ($(A)!0.44!(C)$);
  \draw[given] (E) -- (D);
  \node[vlab] at ($(D)+(0.140,0.143)$)  {$D$};
  \node[vlab] at ($(E)+(-0.162,0.116)$) {$E$};
  \foreach \p/\n/\d in {B/B/below left, C/C/below right, A/A/above}
    {\node[vlab,\d] at (\p) {$\n$};}
\end{tikzpicture}
\figcap{9--33}
\end{bkfigure}}

\newcommand{\FIGIXTHIRTYFOUR}{\begin{bkfigure}[0.54\linewidth]
\begin{tikzpicture}[scale=2.20]
  % 9-34: D at 0.285 of AB, C at the height the right angle at C forces.
  % NO right-angle marks: in Ex. 9-10 #5 both "angle C is right" and
  % "CD perp AB" are what the reader has to PROVE, and the book draws the
  % figure bare (checked on the p.169 plate at 400 dpi).
  \coordinate (A) at (0,0); \coordinate (B) at (2.25,0);
  \coordinate (D) at (0.64,0); \coordinate (C) at (0.64,1.0156);
  \draw[given] (A) -- (B) -- (C) -- cycle;
  \draw[given] (C) -- (D);
  \foreach \p/\n/\d in {A/A/below left, B/B/below right, C/C/above, D/D/below}
    {\node[vlab,\d] at (\p) {$\n$};}
\end{tikzpicture}
\figcap{9--34}
\end{bkfigure}}

\newcommand{\FIGIXTHIRTYFIVE}{\begin{bkfigure}[0.54\linewidth]
\begin{tikzpicture}[scale=2.20]
  % 9-35: same construction with D at 0.712 of AB.  CD perp AB IS given here
  % (Ex. 9-10 *#6), but the book still prints no right-angle mark -- checked
  % on the p.169 plate at 400 dpi -- so none is drawn.
  \coordinate (A) at (0,0); \coordinate (B) at (2.25,0);
  \coordinate (D) at (1.602,0); \coordinate (C) at (1.602,1.0189);
  \draw[given] (A) -- (B) -- (C) -- cycle;
  \draw[given] (C) -- (D);
  \foreach \p/\n/\d in {A/A/below left, B/B/below right, C/C/above, D/D/below}
    {\node[vlab,\d] at (\p) {$\n$};}
\end{tikzpicture}
\figcap{9--35}
\end{bkfigure}}

% --------------------------------------------------------------- Fig 9-36
% A'' on BA and C'' on BC at the same ratio 0.60, matching the 0.60 scale of
% the primed triangle; hence A''C'' || AC.
\newcommand{\FIGIXTHIRTYSIX}{\begin{bkfigure}[0.84\linewidth]
\begin{tikzpicture}[scale=1.667]
  \coordinate (A) at (0,0); \coordinate (C) at (1.65,0);
  \coordinate (B) at (0.85,1.50);
  \draw[given] (A) -- (C) -- (B) -- cycle;
  \coordinate (App) at ($(B)!0.60!(A)$);
  \coordinate (Cpp) at ($(B)!0.60!(C)$);
  \draw[given] (App) -- (Cpp);
  \foreach \p/\n/\d in {A/A/below left, C/C/below right, B/B/above,
                        App/A\IXpr\IXpr/left, Cpp/C\IXpr\IXpr/right}
    {\node[vlab,\d] at (\p) {$\n$};}
  \begin{scope}[xshift=3.05cm, scale=0.60]
    \draw[given] (0,0) -- (1.65,0) -- (0.85,1.50) -- cycle;
    \node[vlab,below left] at (0,0) {$A\IXpr$};
    \node[vlab,below right] at (1.65,0) {$C\IXpr$};
    \node[vlab,above] at (0.85,1.50) {$B\IXpr$};
  \end{scope}
\end{tikzpicture}
\figcap{9--36}
\end{bkfigure}}

% ---------------------------------------------------------- Fig 9-37..9-39
% SPLIT 2026-08-17: Ex. 9-11 cites 9-37 in #2, 9-38 in #3 and 9-39 in #4.
\newcommand{\FIGIXTHIRTYSEVEN}{\begin{bkfigure}[0.58\linewidth]
\begin{tikzpicture}[scale=2.10]
  % 9-37: EA = 10, AC = 14, EB = 15, BD = 21, so EA/EC = EB/ED = 5/12 and
  % AB || CD.  A and B are placed at that ratio, not by eye.
  \coordinate (E) at (0,0); \coordinate (D) at (2.48,0);
  \coordinate (C) at (1.32,1.00);
  \coordinate (A) at ($(E)!5/12!(C)$);
  \coordinate (B) at ($(E)!5/12!(D)$);
  \draw[given] (E) -- (D) -- (C) -- cycle;
  \draw[given] (A) -- (B);
  \node[slab] at ($($(E)!0.5!(A)$)+(-0.133,0.175)$) {10};
  \node[slab] at ($($(A)!0.5!(C)$)+(-0.133,0.175)$) {14};
  \node[slab,below] at ($(E)!0.5!(B)$) {15};
  \node[slab,below] at ($(B)!0.5!(D)$) {21};
  \node[vlab,below left]  at (E) {$E$};  \node[vlab,right] at (D) {$D$};
  \node[vlab,above] at (C) {$C$};
  \node[vlab,above left=1pt] at (A) {$A$}; \node[vlab,below] at (B) {$B$};
\end{tikzpicture}
\figcap{9--37}
\end{bkfigure}}

\newcommand{\FIGIXTHIRTYEIGHT}{\begin{bkfigure}[0.74\linewidth]
% This paired diagram is wider than a single restored-book exercise column.
% Reduce only this exercise figure enough to clear the neighboring column;
% labels and line weights retain the shared book settings.
\begin{tikzpicture}[scale=1.05]
  % 9-38: one hexagon, then the same hexagon rotated -30 and scaled 0.72,
  % so the two really are similar.
  \def\hexA{(-0.87,0.58) (0.41,0.81) (0.79,0.06) (0.39,-0.65)
            (-0.09,-0.62) (-0.64,-0.18)}
  \draw[given] plot coordinates {\hexA} -- cycle;
  \node[vlab,left]  at (-0.87,0.58) {$A$};
  \node[vlab,above] at (0.41,0.81)  {$B$};
  \node[vlab,right] at (0.79,0.06)  {$C$};
  \node[vlab,right] at (0.39,-0.65) {$D$};
  \node[vlab,below] at (-0.09,-0.62){$E$};
  \node[vlab,left]  at (-0.64,-0.18){$F$};
  \begin{scope}[xshift=3.25cm]
    \begin{scope}[rotate=-30, scale=0.72]
      \draw[given] plot coordinates {\hexA} -- cycle;
      \coordinate (ha) at (-0.87,0.58);  \coordinate (hb) at (0.41,0.81);
      \coordinate (hc) at (0.79,0.06);   \coordinate (hd) at (0.39,-0.65);
      \coordinate (he) at (-0.09,-0.62); \coordinate (hf) at (-0.64,-0.18);
    \end{scope}
    \node[vlab,above left]  at (ha) {$A\IXpr$};
    \node[vlab,above right] at (hb) {$B\IXpr$};
    \node[vlab,below right] at (hc) {$C\IXpr$};
    \node[vlab,below]       at (hd) {$D\IXpr$};
    \node[vlab,below left]  at (he) {$E\IXpr$};
    \node[vlab,left]        at (hf) {$F\IXpr$};
  \end{scope}
\end{tikzpicture}
\figcap{9--38}
\end{bkfigure}}

% 9-39: BE _|_ AC with E on AC produced, CD _|_ AB with D the true foot.
% 2026-08-17, against the p.170 plate at 500 dpi: the right-angle square at D
% belongs in the wedge between DB and DC (it was drawn on the A side), and the
% dashed base runs a short way PAST E (it stopped dead at E).
\newcommand{\FIGIXTHIRTYNINE}{\begin{bkfigure}[0.58\linewidth]
\begin{tikzpicture}[scale=2.00]
  \coordinate (A) at (0,0); \coordinate (E) at (2.29,0);
  \coordinate (C) at (1.55,0); \coordinate (B) at (2.29,1.55);
  \draw[given] (A) -- (C) -- (B) -- cycle;
  \coordinate (D) at ($(A)!(C)!(B)$);
  \draw[given] (C) -- (D);
  \draw[aux] (C) -- (2.62,0); \draw[aux] (B) -- (E);
  \draw[fig, line width=0.5pt] (2.29,0.18) -- (2.11,0.18) -- (2.11,0);
  \draw[fig, line width=0.5pt]
    ($(D)!0.18!(B)$) -- ($($(D)!0.18!(B)$)!0.18!(C)$) -- ($(D)!0.18!(C)$);
  \foreach \p/\n/\d in {A/A/below left, B/B/above right, C/C/below,
                        D/D/above left, E/E/below right}
    {\node[vlab,\d] at (\p) {$\n$};}
\end{tikzpicture}
\figcap{9--39}
\end{bkfigure}}

% --------------------------------------------------------- Fig 9-40, 9-41
% Both are body figures (Theorem 9-11 and its proof), so they keep the book's
% side-by-side placement; only the drawing grows.
\newcommand{\FIGIXFORTYFORTYONE}{\begin{bkfigure}[\linewidth]
\begin{minipage}{0.46\linewidth}\centering
\begin{tikzpicture}[scale=1.32]
  % 9-40: right angle at C, D the foot of the altitude (CD^2 = AD.DB).
  \coordinate (A) at (0,0); \coordinate (B) at (2.67,0);
  \coordinate (D) at (1.95,0); \coordinate (C) at (1.95,1.1849);
  \draw[given] (A) -- (B) -- (C) -- cycle;
  \draw[given] (C) -- (D);
  % Both right angles ARE hypotheses of Theorem 9-11, and the book marks both.
  % The foot's box sits on the A side of CD (checked at 400 dpi), and the mark
  % at C is a full square, not an open chevron.
  \draw[rmark] (1.73,0) -- (1.73,0.22) -- (1.95,0.22);
  \draw[rmark] (1.762,1.0706) -- (1.876,0.8826) -- (2.064,0.9969);
  \foreach \p/\n/\d in {A/A/below left, B/B/below right, C/C/above, D/D/below}
    {\node[vlab,\d] at (\p) {$\n$};}
\end{tikzpicture}
\figcap{9--40}
\end{minipage}\hfill
\begin{minipage}{0.44\linewidth}\centering
\begin{tikzpicture}[scale=1.32]
  % 9-41: the same triangle with l', l, l'' through A, D(=C) and B.
  % The apex is labelled C, NOT C' -- Caleb checked the physical copy
  % 2026-08-17 and the p.171 plate at 500 dpi agrees.  The label sits in the
  % wedge between side AC and the dashed line l, just under the apex, and l is
  % named low down beside D exactly as the book names it.
  \coordinate (A) at (0,0); \coordinate (B) at (2.48,0);
  \coordinate (D) at (1.83,0); \coordinate (C) at (1.83,1.0906);
  \draw[given] (A) -- (B) -- (C) -- cycle;
  \draw[aux] (0,-0.38) -- (0,1.62);      \node[vlab,above] at (0,1.62) {$l\IXpr$};
  \draw[aux] (1.83,-0.38) -- (1.83,1.62);
  \draw[aux] (2.48,-0.38) -- (2.48,1.62); \node[vlab,above] at (2.48,1.62) {$l\IXpr\IXpr$};
  \node[vlab] at (1.70,0.26) {$l$};
  \node[vlab] at ($(C)+(-0.17,0.24)$) {$C$};
  \foreach \p/\n/\d in {A/A/below left, B/B/below right, D/D/below right}
    {\node[vlab,\d] at (\p) {$\n$};}
\end{tikzpicture}
\figcap{9--41}
\end{minipage}
\end{bkfigure}}

% ---------------------------------------------------------- Fig 9-42..9-45
% SPLIT 2026-08-17: Ex. 9-12 cites 9-42 in #1, 9-43 in #2, 9-44 in #3/#4 and
% 9-45 in #5, so each stands alone and is set inline after its exercise.
\newcommand{\FIGIXFORTYTWO}{\begin{bkfigure}[0.55\linewidth]
\begin{tikzpicture}[scale=1.90]
  % 9-42: EF = 8, FG = 5, right angle at F; H the foot, so EH : HG = 64 : 25.
  \coordinate (E) at (0,0); \coordinate (G) at (2.6412,0);
  \coordinate (H) at (1.8996,0); \coordinate (F) at (1.8996,1.1873);
  \draw[given] (E) -- (G) -- (F) -- cycle;
  \draw[given] (F) -- (H);
  \node[slab,above left]  at ($(E)!0.5!(F)$) {8};
  \node[slab,above right] at ($(F)!0.5!(G)$) {5};
  \foreach \p/\n/\d in {E/E/below left, G/G/below right, F/F/above, H/H/below}
    {\node[vlab,\d] at (\p) {$\n$};}
\end{tikzpicture}
\figcap{9--42}
\end{bkfigure}}

\newcommand{\FIGIXFORTYTHREE}{\begin{bkfigure}[0.55\linewidth]
\begin{tikzpicture}[scale=1.85]
  % 9-43: a true 5-12-13 right triangle; the altitude on the hypotenuse is
  % 60/13, its foot 25/13 from the short leg.
  \coordinate (A) at (0,0); \coordinate (B) at (2.73,0);
  \coordinate (D) at (0.4038,0); \coordinate (C) at (0.4038,0.9692);
  \draw[given] (A) -- (B) -- (C) -- cycle;
  \draw[given] (C) -- (D);
  \node[slab,left]  at ($(A)!0.5!(C)$) {5};
  \node[slab,above right] at ($(C)!0.5!(B)$) {12};
  \node[slab,below] at ($(A)!0.5!(B)$) {13};
  \node[slab,right] at (0.42,0.50) {$h$};
\end{tikzpicture}
\figcap{9--43}
\end{bkfigure}}

\newcommand{\FIGIXFORTYFOUR}{\begin{bkfigure}[0.55\linewidth]
\begin{tikzpicture}[scale=2.00]
  % 9-44: PS = 4, SR = 5; right angle at Q, QS _|_ PR, so QS = sqrt(20).
  \coordinate (R) at (0,0); \coordinate (P) at (2.52,0);
  \coordinate (S) at (1.40,0); \coordinate (Q) at (1.40,1.2522);
  \draw[given] (R) -- (P) -- (Q) -- cycle;
  \draw[aux] (Q) -- (S);
  \node[slab,below] at ($(R)!0.5!(S)$) {5};
  \node[slab,below] at ($(S)!0.5!(P)$) {4};
  \foreach \p/\n/\d in {R/R/below left, P/P/below right, Q/Q/above, S/S/below}
    {\node[vlab,\d] at (\p) {$\n$};}
\end{tikzpicture}
\figcap{9--44}
\end{bkfigure}}

\newcommand{\FIGIXFORTYFIVE}{\begin{bkfigure}[0.60\linewidth]
\begin{tikzpicture}[scale=1.90]
  % 9-45: right angle at C, right angle at D; b = 2.05, a = 1.50, c = 2.540,
  % so AD = b^2/c and h = ab/c hold exactly.
  \coordinate (A) at (0,0); \coordinate (B) at (2.540,0);
  \coordinate (D) at (1.6545,0); \coordinate (C) at (1.6545,1.2106);
  \draw[given] (A) -- (B) -- (C) -- cycle;
  \draw[aux] (C) -- (D);
  % The book braces ALL THREE sides here (b along AC, a along CB, c along AB)
  % and puts h on the B side of the dashed altitude.
  \draw[bkbrace]  (A) -- (C);
  \draw[bkbrace]  (C) -- (B);
  \draw[bkbracer] (A) -- (B);
  \node[slab] at (0.414,1.170) {$b$};
  \node[slab] at (2.662,1.018) {$a$};
  \node[slab] at (1.92,0.40) {$h$};
  \node[slab,below=3pt] at (1.00,-0.10) {$c$};
  \foreach \p/\n/\d in {A/A/left, B/B/right, C/C/above}
    {\node[vlab,\d] at (\p) {$\n$};}
  \node[vlab] at ($(D)+(-0.30,0.45)$) {$D$};
\end{tikzpicture}
\figcap{9--45}
\end{bkfigure}}

% --------------------------------------------------------------- Fig 9-46
% Right angle at C; D the foot of the altitude, x = AD, y = DB, x.y = h^2.
\newcommand{\FIGIXFORTYSIX}{\begin{bkfigure}[0.58\linewidth]
\begin{tikzpicture}[scale=1.522]
  \coordinate (A) at (0,0); \coordinate (B) at (2.75,0);
  \coordinate (D) at (1.05,0); \coordinate (C) at (1.05,1.3360);
  \draw[given] (A) -- (B) -- (C) -- cycle;
  \draw[aux] (C) -- (D);
  \node[slab,above left]  at ($(A)!0.5!(C)$) {$b$};
  \node[slab,above right] at ($(C)!0.5!(B)$) {$a$};
  \node[slab,above] at (0.44,0) {$x$};
  \node[slab,above] at (1.94,0) {$y$};
  \draw[bkbracer] (A) -- (B);
  \node[slab,below=3pt] at (1.62,-0.10) {$c$};
  \foreach \p/\n/\d in {A/A/left, B/B/right, C/C/above, D/D/above left}
    {\node[vlab,\d] at (\p) {$\n$};}
\end{tikzpicture}
\figcap{9--46}
\end{bkfigure}}

% --------------------------------------------------------------- Fig 9-47
% Two congruent right triangles, legs solid and hypotenuse dashed as printed.
\newcommand{\FIGIXFORTYSEVEN}{\begin{bkfigure}[0.72\linewidth]
\begin{tikzpicture}[scale=1.522]
  \coordinate (C) at (0,0); \coordinate (B) at (0.94,0);
  \coordinate (A) at (0,1.60);
  \draw[given] (C) -- (B); \draw[given] (C) -- (A); \draw[aux] (A) -- (B);
  \node[slab,left]  at ($(C)!0.5!(A)$) {$b$};
  \node[slab,below] at ($(C)!0.5!(B)$) {$a$};
  \node[slab,right] at (0.62,0.85) {$c$};
  \node[vlab,above] at (A) {$A$}; \node[vlab,below left] at (C) {$C$};
  \node[vlab,below right] at (B) {$B$};
  \begin{scope}[xshift=2.35cm]
    \coordinate (P) at (0,0); \coordinate (Q) at (0.94,0);
    \coordinate (R) at (0,1.60);
    \draw[given] (P) -- (Q); \draw[given] (P) -- (R); \draw[aux] (R) -- (Q);
    \node[slab,left]  at ($(P)!0.5!(R)$) {$b$};
    \node[slab,below] at ($(P)!0.5!(Q)$) {$a$};
    \node[vlab,above] at (R) {$R$}; \node[vlab,below left] at (P) {$P$};
    \node[vlab,below right] at (Q) {$Q$};
  \end{scope}
\end{tikzpicture}
\figcap{9--47}
\end{bkfigure}}

% --------------------------------------------------------- Fig 9-48, 9-49
% SPLIT 2026-08-17: Ex. 9-13 cites 9-48 in *#4 and 9-49 in #8.
%
% 9-48: SCHEMATIC, by Caleb's decision of 2026-08-17.  The earlier draft drew
% the diagonals at the exercise's own BD : AC = 24 : 10, which hands the reader
% the answer; the book draws an ordinary-looking rhombus at about 1.58 : 1 and
% lets the numbers do the work.  The three hypotheses the exercise DOES state
% -- AB || CD, BC || AD, AC _|_ BD -- still hold exactly: the figure is built
% from two perpendicular half-diagonals, which forces all four sides equal.
\newcommand{\FIGIXFORTYEIGHT}{\begin{bkfigure}[0.52\linewidth]
\begin{tikzpicture}[scale=3.20]
  % u = (0.450,-0.893) along AC ; v = (0.893,0.450) along BD ; |AC|/2 = 0.52,
  % |BD|/2 = 0.82, matching the book's proportions on the p.174 plate.
  \coordinate (A) at (-0.234, 0.464);  \coordinate (C) at (0.234,-0.464);
  \coordinate (B) at (-0.732,-0.369);  \coordinate (D) at (0.732, 0.369);
  \draw[given] (A) -- (D) -- (C) -- (B) -- cycle;
  \draw[given] (A) -- (C); \draw[given] (B) -- (D);
  \foreach \p/\n/\d in {A/A/above left, D/D/above right,
                        B/B/below left, C/C/below right}
    {\node[vlab,\d] at (\p) {$\n$};}
\end{tikzpicture}
\figcap{9--48}
\end{bkfigure}}

\newcommand{\FIGIXFORTYNINE}{\begin{bkfigure}[0.56\linewidth]
\begin{tikzpicture}[scale=2.20]
  % 9-49: OA = 21 north, OB = 20 east, right angle at O; AB dashed across the
  % lake.  The four contours are traced from the iPad scan (scan11 p8).
  \coordinate (O) at (0,0); \coordinate (A) at (0,2.10);
  \coordinate (B) at (2.00,0);
  \draw[given, line width=0.5pt]
    plot[smooth cycle, tension=0.7] coordinates {
      (1.025,2.024) (1.424,1.943) (1.833,1.736) (2.132,1.504) (2.309,1.246)
      (2.324,1.029) (2.242,0.781) (2.051,0.562) (1.804,0.452) (1.560,0.397)
      (1.312,0.372) (1.066,0.375) (0.848,0.431) (0.671,0.522) (0.520,0.659)
      (0.398,0.829) (0.306,1.029) (0.276,1.224) (0.306,1.422) (0.368,1.599)
      (0.464,1.777) (0.645,1.928) (0.823,2.002)};
  \draw[given, line width=0.5pt]
    plot[smooth cycle, tension=0.7] coordinates {
      (1.015,1.847) (1.339,1.777) (1.627,1.629) (1.859,1.449) (2.025,1.271)
      (2.107,1.079) (2.092,0.892) (1.996,0.740) (1.804,0.648) (1.586,0.603)
      (1.368,0.589) (1.150,0.611) (0.944,0.670) (0.767,0.765) (0.630,0.902)
      (0.528,1.068) (0.479,1.261) (0.516,1.437) (0.601,1.599) (0.782,1.750)};
  \draw[given, line width=0.5pt]
    plot[smooth cycle, tension=0.7] coordinates {
      (1.025,1.573) (1.312,1.496) (1.560,1.353) (1.738,1.191) (1.833,1.029)
      (1.812,0.877) (1.697,0.765) (1.519,0.710) (1.320,0.703) (1.136,0.748)
      (0.985,0.829) (0.884,0.950) (0.845,1.101) (0.870,1.267) (0.929,1.422)};
  \draw[given, line width=0.5pt]
    plot[smooth cycle, tension=0.7] coordinates {
      (1.150,1.271) (1.310,1.304) (1.428,1.267) (1.482,1.193) (1.457,1.120)
      (1.361,1.075) (1.247,1.056) (1.150,1.095) (1.099,1.164) (1.111,1.230)};
  \draw[given] (A) -- (O) -- (B);
  \draw[aux] (A) -- (B);
  \node[vlab,above] at (A) {$A$};
  \node[vlab,below left] at (O) {$O$};
  \node[vlab,right] at (B) {$B$};
\end{tikzpicture}
\figcap{9--49}
\end{bkfigure}}

% --------------------------------------------------------- Fig 9-50, 9-51
\newcommand{\FIGIXFIFTYFIFTYONE}{\begin{bkfigure}[\linewidth]
\begin{minipage}{0.38\linewidth}\centering
\begin{tikzpicture}[scale=1.45]
  % 9-50: C is the midpoint of AB, P at height sqrt(r^2 - (AB/2)^2); the two
  % dashed arcs are the circles of radius r about A and B.
  % REBUILT 2026-08-17 from the book photograph (feedback/small/20260817_102830):
  %   * the apex was at 0.65 of AB, the book's at 0.50, which made the lens far
  %     too fat (arcs cutting AB at 18%/82% instead of the book's 31%/69%) and
  %     drove three thick strokes into a blob at P that read as an arrowhead;
  %   * the arcs stopped dead at their intersections.  The book OVERSHOOTS each
  %     arc so a small X is struck at P and again at the lower crossing -- the
  %     distinctive mark of the construction.  Each arc now runs to +-57 deg,
  %     i.e. 12 deg past the +-45 deg where the two meet;
  %   * the arcs are drawn finer than the triangle, as in the book, and PC is
  %     no longer heavier than PA and PB.
  \coordinate (A) at (0,0); \coordinate (B) at (1.80,0);
  \coordinate (C) at (0.90,0); \coordinate (P) at (0.90,0.90);
  \draw[finedash] ([shift=(-57:1.27279)]A) arc (-57:57:1.27279);
  \draw[finedash] ([shift=(123:1.27279)]B) arc (123:237:1.27279);
  \draw[given] (A) -- (B);
  \draw[given] (C) -- ($(C)!0.985!(P)$);
  \draw[given] (A) -- (P) -- (B);
  \node[slab] at (0.337,0.563) {$r$};
  \foreach \p/\n/\d in {A/A/below left, B/B/below right, C/C/below}
    {\node[vlab,\d] at (\p) {$\n$};}
  \node[vlab] at (0.90,1.15) {$P$};   % in the gap between the two arc tails
\end{tikzpicture}
\figcap{9--50}
\end{minipage}\hfill
\begin{minipage}{0.55\linewidth}\centering
\begin{tikzpicture}[scale=1.60]
  % 9-51: C on the semicircle, D the foot of the perpendicular from C, with
  % x = AD = r - y, y = DO, z = CD = sqrt(r^2 - y^2), r the radius.
  \def\R{1.40}
  \coordinate (O) at (0,0);
  \coordinate (A) at (-\R,0); \coordinate (B) at (\R,0);
  \coordinate (D) at (-0.637,0); \coordinate (C) at (-0.637,1.2464);
  \draw[given] (A) -- (B);
  \draw[given] (A) arc (180:0:\R);
  \draw[given] (A) -- (C) -- (B);
  \draw[given] (C) -- (D);
  \draw[given] (C) -- (O);
  \dt{O}
  \node[vlab,below left]  at (A) {$A$};
  \node[vlab,below right] at (B) {$B$};
  \node[vlab,above] at (C) {$C$};
  \node[vlab,below] at (D) {$D$};
  \node[vlab,below] at (O) {$O$};
  \node[slab,above=1pt] at (-1.02,0) {$x$};
  \node[slab,above=1pt] at (-0.32,0) {$y$};
  \node[slab] at (-0.80,0.34) {$z$};
  \node[slab,right=1pt] at (-0.3185,0.6232) {$r$};
  \node[slab,above=1pt] at (0.70,0) {$r$};
\end{tikzpicture}
\figcap{9--51}
\end{minipage}
\end{bkfigure}}

% --------------------------------------------------------------- Fig 9-52
% AB a vertical diameter; CD _|_ AB at E, with EO = 12 and radius 15, so the
% half-chord is 9.
\newcommand{\FIGIXFIFTYTWO}{\begin{bkfigure}[0.56\linewidth]
\begin{tikzpicture}[scale=1.75]
  \draw[given] (0,0) circle (1.5);
  \coordinate (A) at (0,1.5); \coordinate (B) at (0,-1.5);
  \coordinate (E) at (0,1.2);
  \coordinate (C) at (-0.9,1.2); \coordinate (D) at (0.9,1.2);
  \draw[given] (A) -- (B);
  \draw[given] (C) -- (D);
  \dt{0,0}
  \node[vlab,above] at (A) {$A$};   \node[vlab,below] at (B) {$B$};
  \node[vlab] at ($(C)+(-0.30,0.18)$) {$C$};
  \node[vlab] at ($(D)+(0.30,0.18)$) {$D$};
  \node[vlab,below right] at (E) {$E$};
  \node[vlab,right] at (0,0) {$O$};
\end{tikzpicture}
\figcap{9--52}
\end{bkfigure}}

% ------------------------------------------ Figs 9-53/9-54, 9-55 (hoisted
% from ch09.tex local block after the Review-Exercise transcription pass)
% Figures 9-53, 9-54, 9-55 belong to the Review Exercises (book pp.178-179) and
% did not exist when figures09.tex was reviewed; that file is owned by the
% figure-review pass, so they are defined here per the "new local macro" rule.
% HOIST: move these three into figures09.tex and add their hypotheses to
% tools/constraints/ch09.py.  See chapters/ch09-UNCERTAIN.md sec. 3.
%
% 9-53: right trapezoid, bases 18 (bottom) and 12, altitude 8; the slant leg is
% 10, so its run is 6 (6-8-10), which is exactly the 18-12 difference.  The
% nonparallel sides extended meet at (18,24): the enclosing triangle is 18-24-30
% with the right angle on the vertical side, which is what Ex. 11 asks students
% to show.  Every vertex is forced by those numbers; nothing is placed by eye.
% SPLIT 2026-08-17: Review Ex. 11 cites 9-53, *Ex. 15 cites 9-54 and Ex. 18
% cites 9-55, so all three are set inline after their own exercise.
\newcommand{\FIGIXFIFTYTHREE}{\begin{bkfigure}[0.46\linewidth]
\begin{tikzpicture}[x=0.165cm,y=0.165cm]
  \coordinate (A) at (0,0);   \coordinate (B) at (18,0);
  \coordinate (P) at (18,24); \coordinate (D) at (6,8);
  \coordinate (C) at (18,8);  \coordinate (H) at (6,0);
  \draw[given] (A) -- (B) -- (P) -- cycle;
  \draw[given] (D) -- (C);
  \draw[given] (D) -- (H);
  \node[slab,left]  at ($($(A)!0.5!(D)$)+(-0.9,0.7)$) {10};
  \node[slab,right] at ($(D)!0.5!(H)$) {8};
  \node[slab,above] at ($(D)!0.5!(C)$) {12};
  \node[slab,below] at ($(A)!0.5!(B)$) {18};
\end{tikzpicture}
\figcap{9--53}
\end{bkfigure}}

% 9-54: Menelaus.  D is fixed at 2/3 of AB and F on ray AC beyond C; E is then
% the true intersection of DF with BC, so AD/DB . BE/EC . CF/FA = 1 holds
% exactly by construction.  X, Y, Z are projection feet, never eyeballed.
% REBUILT 2026-08-17 from the book photograph (feedback/small/20260817_102917).
% The render's triangle was tall and its apex central (apex 53% along AC,
% height 0.89 of the base); the book's is flat and wide with the apex about
% 39% along and 0.67 of the base high.  That mattered: at the old shape the
% dashed AX ran within 5 deg of side AB and the two were indistinguishable for
% a third of their length.  At the book's shape they diverge by 15 deg.  The
% base also rises slightly to the right, as the book's does, the l/X/D labels
% are pulled apart along and across the line, Y and Z now hug their own feet,
% and the dashed extension of AC uses a longer dash so it does not read as a
% fourth perpendicular.
\newcommand{\FIGIXFIFTYFOUR}{\begin{bkfigure}[0.92\linewidth]
\begin{tikzpicture}[scale=1.28]
  \coordinate (A) at (0,0); \coordinate (B) at (1.24,2.21);
  \coordinate (C) at (3.18,0.20);
  \coordinate (F) at ($(A)!1.70!(C)$);
  \coordinate (D) at ($(A)!2/3!(B)$);
  \coordinate (E) at (intersection of D--F and B--C);
  \coordinate (L) at ($(D)!-0.30!(F)$);
  \coordinate (G) at ($(D)!1.10!(F)$);
  \coordinate (X) at ($(D)!(A)!(F)$);
  \coordinate (Y) at ($(D)!(B)!(F)$);
  \coordinate (Z) at ($(D)!(C)!(F)$);
  \draw[given] (A) -- (B) -- (C) -- cycle;
  \draw[longdash] (C) -- ($(A)!1.80!(C)$);
  \draw[given] (L) -- (G);
  \draw[aux]   (A) -- (X);
  \draw[aux]   (B) -- (Y);
  \draw[aux]   (C) -- (Z);
  % Every letter is offset along l's own normal n = (0.2402, 0.9707): the ones
  % the book sets above the line get +n, the ones below get -n, so none of them
  % lands on l or on the perpendicular that meets it there.
  \node[vlab] at ($(L)+(-0.06,0.20)$) {$l$};
  % X and D printed as one blob "XD" in the render Caleb photographed, so they
  % are pushed apart ALONG l as well as across it: 0.43 units centre to centre.
  \node[vlab] at ($(X)+(-0.044,0.238)$) {$X$};
  \node[vlab] at ($(D)+(-0.063,0.384)$) {$D$};
  \node[vlab] at ($(Y)+(-0.058,-0.233)$) {$Y$};
  \node[vlab] at ($(E)+(0.058,0.233)$)  {$E$};
  \node[vlab] at ($(Z)+(0.053,0.214)$) {$Z$};
  \node[vlab,above]      at (B) {$B$};
  \node[vlab,below left] at (A) {$A$};
  \node[vlab,below]      at (C) {$C$};
  \node[vlab,below]      at (F) {$F$};
\end{tikzpicture}
\figcap{9--54}
\end{bkfigure}}

% 9-55: D and E are the feet of the altitudes from C and from B, placed with the
% projection operator, so the two right angles are exact.
% 2026-08-17, against the p.179 plate at 500 dpi: the right-angle square at E
% was struck at 0.16 of |EA| and 0.16 of |EB| -- i.e. proportional to sides
% four times longer than the ones at D -- so it printed far larger than D's and
% hung into the triangle.  Both squares are now 0.17 units on a side, as the
% book draws them, and E's sits in the wedge between EC and EB (the book's
% side).  The triangle is also flattened from 0.85 to the book's 0.72.
\newcommand{\FIGIXFIFTYFIVE}{\begin{bkfigure}[0.72\linewidth]
\begin{tikzpicture}[scale=1.90]
  \coordinate (A) at (0,0); \coordinate (B) at (3.43,0);
  \coordinate (C) at (2.05,2.46);
  \coordinate (D) at ($(A)!(C)!(B)$);
  \coordinate (E) at ($(A)!(B)!(C)$);
  \draw[given] (A) -- (B) -- (C) -- cycle;
  \draw[given] (C) -- (D);
  \draw[given] (B) -- (E);
  % D = (2.05, 0); E = (1.40575, 1.68689) with unit vectors E->C (0.64018,
  % 0.76821) and E->B (0.76821, -0.64019).  Both squares are 0.17 a side.
  \draw[tick] (1.88,0) -- (1.88,0.17) -- (2.05,0.17);
  \draw[tick] (1.51458,1.81749) -- (1.64518,1.70866) -- (1.53635,1.57806);
  \node[vlab,above]       at (C) {$C$};
  \node[vlab,below left]  at (A) {$A$};
  \node[vlab,below right] at (B) {$B$};
  \node[vlab,below]       at (D) {$D$};
  \node[vlab,above left]  at (E) {$E$};
\end{tikzpicture}
\figcap{9--55}
\end{bkfigure}}
